\chapter{Campionamento}

\begin{figure}[h]
    \centering
    \includegraphics[scale = 0.2]{Campionamento1.png}
\end{figure}  

\newpage 

\section{Campionamento di segnale analogico}

\footnote{ 
    Slide del prof | 4 - Campionamento | pag 2\\
    Appunti| 2025-10-14 | pag 1
}

Per campionamento si intende quel processo che, ad ogni periodo di campionamento T detto periodo di campionamento, 
si rivela che valore ha il segnale analogico in quegli istanti. \newline 

Il reciproco di T periodo di campionamento è chiamato $F_s$ frequenza di campionamento. \newline 

In formule: 

{
    \Large 
    \begin{equation}
        F_s = \dfrac{1}{T}
    \end{equation}
}

Da un punto di vista del sistema, nel campionamento avremo questo sistema: 

\begin{figure}[h]
    \centering
    \includegraphics[scale = 0.4]{Sistema campionamento ideale.png}
\end{figure}  

dove, visualizzando la figura da sinistra verso destra, avremo: 

\begin{itemize}
    \item $x_c \left(t\right)$ è il segnale analogico da campionare (e che quindi ha valori continui sia nell'asse delle ampiezze che nel tempo) 
    \item $s \left(t\right)$ è una sequenza di impulsi di Dirac continuo 
    \item $x_s \left(t\right)$ è il segnale campionato
\end{itemize}

\begin{tcolorbox}[colback=red!5!white,colframe=red!75!black] \label{eq: 1}
Quale è la relazione tra $x_s \left( t\right)$ e $x \left[ n\right]$ ? \newline 

È $x_s \left( t\right)$ o $x \left[ n\right]$ il segnale quantizzato? \newline 

Il convertitore è un sistema LTI che fa la convoluzione tra chi? \newline 

(Vedi anche grafico dell'andamento dei segnali dentro la slide)
\end{tcolorbox}

Analizzando la sequenza $s \left(t\right)$ avremo che: 

{
    \Large 
    \begin{equation}
        s \left( t\right) 
        = 
        \sum_{n = - \infty}^{+ \infty} 
        \delta \left(t - n \cdot T\right)
    \end{equation}
} 

dove $\delta \left(t - n \cdot T\right)$ è una delta di Dirac di impulso continuo che si ripete ogni multiplo di periodo T. \newline 

Allora possiamo scrivere: 

{
    \Large 
    \begin{equation}
        \begin{split}
            x_s \left(t\right) 
            &= 
            x_c \left(t\right)  \cdot s \left(t\right) 
            \\
            &=  
            x_c \left(t\right)  
            \cdot 
            \sum_{n = - \infty}^{+ \infty} 
            \delta \left(t - n \cdot T\right) 
        \end{split}
    \end{equation}
}

Grazie alla proprietà dell'impulso, 
possiamo portare il segnale analogico $x_c \left(t\right)$ dentro la sommatoria, quindi: 

{
    \Large 
    \begin{equation}
        \begin{split}
            x_s \left(t\right) 
            &= 
            x_c \left(t\right)  
            \cdot 
            \sum_{n = - \infty}^{+ \infty} 
            \delta \left(t - n \cdot T\right)  
            \\
            &\downarrow
            \\
            x_s \left(t\right) 
            &= 
            \sum_{n = - \infty}^{+ \infty} 
            x_c \left(n \cdot T\right)  
            \cdot 
            \delta \left(t - n \cdot T\right)  
        \end{split}
    \end{equation}
}

Visualizzando l'andamento del segnale, avremo il seguente caso: 

\begin{figure}[h]
    \centering
    \includegraphics[scale = 0.25]{Da segnale analogico a segnale campionato.png}
\end{figure}  

\newpage 

\subsection{Campionamento: analisi in frequenza}

\footnote{ 
    Slide del prof | 4 - Campionamento | pag 3\\
    Appunti| 2025-10-14 | pag 1
}

Per le proprietà della trasformata di Fourier continua, possiamo scrivere che: 

{
    \Large 
    \begin{equation}
        F \left[ x_c \left(t\right) \cdot s \left(t\right)\right] 
        = 
        \dfrac{1}{2 \pi} 
        X_c \left(\jmath \Omega \right) * S \left(\jmath \Omega \right)
    \end{equation}
}

cioè il prodotto tra $x_c \left(t\right)$ e $s \left(t\right)$ nel tempo è la convoluzione in frequenza di $X_c \left(\jmath \Omega \right)$ e $S \left(\jmath \Omega \right)$, 
che sono, rispettivamente, le trasformate di Fourier dei segnali $x_c \left(t\right)$ e $s \left(t\right)$. \newline 

Si noti che le trasformate sono a tempo continuo e dipendente dalla frequenza continua (non periodica). \newline 

\begin{tcolorbox}[colback=red!5!white,colframe=red!75!black] \label{eq: 2} 
Non ho capito l'ultima frase \newline 

"Si noti che le trasformate sono a tempo continuo e dipendente dalla frequenza continua (non periodica)"

\end{tcolorbox}

Poichè risulta: 

{
    \Large 
    \begin{equation}
        F \left[ s \left(t\right)\right]
        = 
        \dfrac{2 \pi}{T}
        \sum_{k = -\infty}^{+ \infty} 
        \delta\left(\Omega - k \Omega_s\right)
    \end{equation}
} 

\begin{tcolorbox}[colback=red!5!white,colframe=red!75!black] \label{eq: 3} 
La trasformata di una Delta di Dirac nel tempo non era uno spettro piatto in frequenza? \newline 

Perchè abbiamo scritto che la trasformata di Fourier di una Delta di Dirac equispaziati nel tempo è a sua volta una sequenza di Delta di Dirac nel tempo? \newline 

Perchè nel campionamento il tempo è discretizzato e non continuo? 
\end{tcolorbox}

dove $\Omega_s$ prende il nome di pulsazione di campionamento perchè: 

{
    \Large 
    \begin{equation}
        \Omega_s 
        = 
        \dfrac{2 \pi}{T}
    \end{equation}
}

dove, a sua volta, T è il periodo di campionamento. \newline 

\newpage 

Allora possiamo scrivere che, facendo tutti i passaggi e le dovute sostituzioni: 

{
    \Large 
    \begin{equation}
        \begin{split}
        F \left[ x_c \left(t\right) \cdot s \left(t\right)\right] 
        &= 
        F \left[ x_s \left( t\right)\right] 
        \\
        &= 
        X_s \left(\jmath \Omega\right) 
        \\
        &= 
        \dfrac{1}{2 \pi} 
        X_c \left(\jmath \Omega \right) * S \left(\jmath \Omega \right) 
        \\
        &=
        \dfrac{1}{2 \pi} 
        X_c \left(\jmath \Omega \right) 
        * 
        \left[
        \dfrac{2 \pi}{T}
        \sum_{k = -\infty}^{+ \infty} 
        \delta\left(\Omega - k \Omega_s\right)
        \right]
        \\
        &= 
        \dfrac{1}{T}
        X_c \left(\jmath \Omega \right) 
        *
        \sum_{k = -\infty}^{+ \infty} 
        \delta\left(\Omega - k \Omega_s\right) 
        \\
        &= 
        \dfrac{1}{T}
        \sum_{k = -\infty}^{+ \infty} 
        X_c \left(\jmath \Omega - k \jmath \Omega_s\right) 
        \end{split}
    \end{equation}
} 

dove la convoluzione della trasformata di Fourier di $x_c \left( t\right)$, cioè $X_c \left(j \Omega\right)$, 
e la sequenza di Delte di Dirac equispaziate in frequenza di multipli della frequenza di campionamento $\Omega_s$ 
è uguale allo spettro di $X_c (\jmath)$ equispaziato ogni multiplo di $\jmath \Omega_s$. \newline 

Riassumendo: 

{
    \Large 
    \begin{equation}
    X_s \left(\jmath \Omega\right)  
    =    
    \dfrac{1}{T}
    \sum_{k = -\infty}^{+ \infty} 
    X_c \left(\jmath \Omega - k \jmath \Omega_s\right)
    \end{equation}
}

\newpage 

\subsection{Campionamento e ricostruzione}

\footnote{ 
    Slide del prof | 4 - Campionamento | pag 4\\
    Appunti| 2025-10-14 | pag 1
}

Ricordando che, per il campionamento abbiamo questo schema: 

\begin{figure}[h]
    \centering
    \includegraphics[scale = 0.25]{Sistema campionamento ideale.png}
\end{figure}  

analizziamo in Fourier i vari segnali indicati in figura nello spettro. \newline 

Se consideriamo $x_c \left( t\right)$ segnale analogico da campionare, 
abbiamo che, in Fourier, avrà il sequente spettro (di esempio): 

\begin{figure}[h]
    \centering
    \includegraphics[scale = 0.3]{Spettro in pulsazioni di un segnale analogico di esempio.png}
\end{figure}  

Come si nota dalla figura, lo spettro di $X_c \left(\jmath \Omega\right)$ si estende nell'intervallo delle pulsazioni $\left[ - \Omega_N, + \Omega\right]$. \newline 

Se la pulsazione di campionamento $\Omega_s$ è oppurtunamente scelta, 
lo spettro in Fourier del segnale campionato $x_s \left(t\right)$, cioè $X_s \left(\jmath \Omega\right)$, 
avrà il seguente spettro: 

\begin{figure}[h]
    \centering
    \includegraphics[scale = 0.3]{Esempio dello spettro di un segnale campionato.png}
\end{figure}  

dove in blu è indicato l'andamento di $X_s \left(\jmath \Omega\right)$. \newline 

Una osservazione che possimo svolgere sullo spettro del segnale campionato $X_s \left(\jmath \Omega\right)$ 
è che lo spettro del segnale si ripete. \newline  

\newpage 

Questo avviene perchè la discretizzazione nel tempo del segnale, come si può notare nella seguente figura: 

\begin{figure}[h]
    \centering
    \includegraphics[scale = 0.3]{Esempio segnale campionato.png}
\end{figure}  

crea replicità in frequenza. \newline 

\begin{tcolorbox}
    Invece nel corso di Teoria di Segnali avevamo visto il contrario: se il segnale è periodico nel tempo, allora avremo discretizzazione nello spettro, 
    ma, come ben sappiamo, grazie alle proprietà tra Fourier e il tempo, se vale un principio in un dominio vale anche nell'altro e viceversa.
\end{tcolorbox}

\begin{tcolorbox}[colback=green!5!white,colframe=green!75!black]
TO-DO: Inserire seconda figura (quella con su scritta Aliasing) e scriverci le osservazioni 
\end{tcolorbox} 

\begin{tcolorbox}[colback=red!5!white,colframe=red!75!black] \label{eq: 4}
Il commento nella slide: \newline 

"Questa modulazione provoca aliasing che deve essere cancellato opportunamente" \newline 

è riferito alla seconda figura partendo dall'alto della slide? \newline 

In quella figura, si è considerato un segnale sotto-campionato, ecco perchè esiste aliasing?
\end{tcolorbox} 

Ritornando alla figura dello spettro di $X_s \left( \jmath \Omega \right)$: 

\begin{figure}[h]
    \centering
    \includegraphics[scale = 0.3]{Esempio dello spettro di un segnale campionato.png}
\end{figure}  

Considerando una finestra rettangolare (indicata in figura in colore rosso) nell'intervallo 
$\left[ - \Omega_N, + \Omega_N\right]$, 
possiamo riottenere lo spettro del segnale originale, proprio come si vede nella seguente figura:   

\begin{figure}[h]
    \centering
    \includegraphics[scale = 0.3]{Spettro del segnale ricostruito esempio.png}
\end{figure}  

dove, lo spettro di $X_R \left(\jmath \Omega\right)$ è uguale allo spettro del segnale originale $X_c \left(\jmath \Omega\right)$. \newline 

\newpage 

\section{Teorema di Nyquist} \label{sec: Teorema di Nyquist}

\footnote{ 
    Slide del prof | 4 - Campionamento | pag 5\\
    Appunti| 2025-10-14 | pag 1
}

\begin{tcolorbox}
    Come studiato nei corsi precedenti, anche in questo corso, è molto utile il Teorema di Nyquist. 
\end{tcolorbox} 

Nelle pagine precedenti, si è potuto ricostruire il segnale originale da quello campionato perchè il campionamrnto ha soddisfatto il Teorema di Nyquist. \newline 

Di seguito il teorema di Nyquist. \newline 

Se il segnale da campionare $x_c \left(t\right)$ è limitato in banda, cioè: 

{
    \Large 
    \begin{equation}
        X_c \left(\jmath \Omega\right) = 0 \text{ per } \abs{\Omega} > \Omega_N
    \end{equation}
}

allora $x_c \left(t\right)$ è univocamente rappresentato dai suoi campioni: 

{
    \Large 
    \begin{equation}
        x \left[n \right]
        = 
        x_c \left(n \cdot T\right) 
        \text{ per }
        n = 0, \pm 1, \pm 2, \dots
    \end{equation}
}

se risulta: 

{
    \Large 
    \begin{equation}
        \Omega_s = \dfrac{2 \pi}{T} 
        > 
        2 \Omega_N
    \end{equation}
}

dove: 

\begin{itemize}
    \item T è il periodo di campionamento 
    \item $\Omega_s$ è la pulsazione di campionamento 
\end{itemize} 

$\Omega_N$ prende il nome di pulsazione di Nyquist. \newline 

$\blacksquare $ \newline 

Le grandezze di nostro interesse per questo corso sono le seguenti: 

\begin{itemize}
    \item $f_{s}$ è la frequenza di campionamento di Nyquist, che vale: 
    {
        \Large 
        \begin{equation}
            f_s = \dfrac{\Omega_s}{2 \pi}
        \end{equation}
    }
    \item $f_{Ny}$, cioè la frequenza di Nyquist, che vale: 
    {
        \Large 
        \begin{equation}
            f_{Ny} = \dfrac{\dfrac{\Omega_s}{2}}{2 \pi}
        \end{equation}
    }
    
\end{itemize} 

\newpage 

\section{Relazione tra spettri continui e discreti}

\footnote{ 
    Slide del prof | 4 - Campionamento | pag 6\\
    Appunti| 2025-10-14 | pag 1
}   

Consideriamo dal punto di vista matematico la relazione tra 
un segnale a spettro continuo $x_s \left(t\right)$ 
e il segnale a spettro discreto $x \left[ n\right]$. \newline 

Sapendo che nel dominio del tempo t, il segnale campionato $x_s \left(t\right)$ la possiamo scrivere come: 

{
    \Large 
    \begin{equation}
        x_s \left(t\right) 
        = 
        \sum_{n = - \infty}^{+ \infty}
        x_c \left(n T\right) \cdot \delta \left( t - nT\right)
    \end{equation}
} 

Allora la sua trasformata di Fourier è la seguente: 

{
    \Large 
    \begin{equation}
        \begin{split}
        x_s \left(t\right) 
        &= 
        \sum_{n = - \infty}^{+ \infty}
        x_c \left(n T\right) \cdot \delta \left( t - nT\right)
        \\
        &\updownarrow
        \\
        X_s \left(\jmath \Omega\right)
        &= 
        \sum_{n = -\infty}^{+ \infty}
        x_c \left(n T\right)
        \cdot 
        e^{-\jmath \Omega T n}
        \end{split}
    \end{equation}
}

Ma sapendo che il segnale discreto $x \left[ n\right]$ in Fourier è: 

{
    \Large 
    \begin{equation}
        X \left(e^{\jmath \omega}\right) 
        = 
        \sum_{n = -\infty}^{+ \infty}
        x \left[n\right] \cdot e^{- \jmath \omega n}
    \end{equation}
}

Allora possiamo dire che sia il segnale campionato $x_s \left(t\right)$ che il segnale campionato a tempo discreto $x \left[ n\right]$ sono uguali se: 

{
    \Large 
    \begin{equation}
        X_s \left(\jmath \Omega\right) = X \left(e^{\jmath \omega}\right)
        \text{ se }
        \omega = \Omega T
    \end{equation}
}

Scritta in parole, se lo spettro continuo del segnale campionato $X_s \left(\jmath \Omega\right)$ viene valutato ogni $\Omega T$, 
allora lo spettro continuo $X_s \left(\jmath \Omega\right)$ e lo spettro del segnale a tempo discreto $X \left(e^{\jmath \omega}\right)$ sono uguali. \newline 

In particolare: 

{
    \Large 
    \begin{equation}
        X_s \left(\jmath \Omega \right)
        = 
        \left. 
        X \left(e^{\jmath \omega}\right)
        \right|_{\omega = \Omega T} 
        = 
        \dfrac{1}{T}
        \sum_{k = -\infty}^{+ \infty}
        X_c \left(\jmath \Omega - \jmath \dfrac{2 \pi k}{T}\right)
    \end{equation}
}

\newpage 

Grazie a queste relazioni, possiamo visualizzare queste immagini di esempio: 

\begin{figure}[h]
    \centering
    \includegraphics[scale = 0.4]{Esempio di spettri continui e discreti dello stesso segnale.png}
\end{figure}  

Notiamo che gli spettri dei segnali (le due figure a destra) sono identitici: 
ciò che cambia è solo l'asse delle ascisse che al posto di avere le pulsazione in $\Omega$ avremo le pulsazioni in $\omega$. \newline 

\begin{tcolorbox}[colback=red!5!white,colframe=red!75!black] \label{eq: 5}
Leggi ultima frase scritta. \newline 

Quindi come chiamamo $\omega$, pulsazione nel dominio discreto? 
\end{tcolorbox} 

\newpage 

\subsection{Ricostruzione del segnale analogico}

\footnote{ 
    Slide del prof | 4 - Campionamento | pag 7 - 8\\
    Appunti| 2025-10-14 | pag 1
}   

Ricordando che, per il campionamento abbiamo questo schema: 

\begin{figure}[h]
    \centering
    \includegraphics[scale = 0.2]{Sistema campionamento ideale.png}
\end{figure}  

e lo spettro dei segnali: 

\begin{figure}[h]
    \centering
    \includegraphics[scale = 0.25]{Esempio di spettri continui e discreti dello stesso segnale.png}
\end{figure}  

allora possiamo scrivere che il segnale a $x \left[ n\right]$ ha il seguente andamento: 

{
    \Large 
    \begin{equation}
        x \left[ n\right]
        = 
        x_s \left(t\right)
        = 
        \sum_{n = -\infty}^{+ \infty}
        x \left[ n\right]
        \cdot 
        \delta \left(t - nT\right)
    \end{equation}
}

Se non c'è aliasing (cioè si rispetta il \ref{sec: Teorema di Nyquist} \nameref{sec: Teorema di Nyquist}), 
allora si può ricostruire il segnale continuo $x_r \left(t\right)$ mediante un filtro passa-basso continuo ideale con la seguente funzione di trasferimento $H_r \left(\jmath \Omega\right)$: 

{
    \Large 
    \begin{equation}
        H_r \left(\jmath \Omega\right) 
        = 
        \begin{cases}
            T \text{ per } \abs{\Omega } < \dfrac{\Omega_s}{2} 
            \\
            0 \text{ altrove}
        \end{cases}
    \end{equation}
}

dove: 

\begin{itemize}
    \item T è il guadagno del filtro (e non il periodo del segnale)
    \item $\Omega_s$ è la pulsazione di taglio, chiamata anche pulsazione di cutoff
\end{itemize}

Da un punto di vista grafico, $H_r \left(\jmath \right)$ avrà questo andamento nello spettro: 

\begin{figure}[h]
    \centering
    \includegraphics[scale = 0.4]{Filtro interpolatore in frequenza.png}
\end{figure}  

Questo filtro prende il nome di filtro interpolatore, cioè il filtro interpolatore è una finestra in frequenza. \newline 

Sapendo che la trasformata di un filtro continuo in frequenza è un sinc continuo, 
allora avremo che nel tempo il filtro interpolatore avrà questo andamento: 

{
    \Large 
    \begin{equation}
        h_r \left(t\right)
        = 
        \dfrac{\sin \left(\dfrac{\Omega_s}{2} t\right)}{\dfrac{\Omega_s}{2} t}
    \end{equation}
}

che, dal punto di vista grafico: 

\begin{figure}[h]
    \centering
    \includegraphics[scale = 0.4]{Filtro interpolatore nel tempo.png}
\end{figure}  

Essendo un sinc continuo, si estende nel tempo nell'intervallo $\left[ - \infty, + \infty\right]$. \newline 

Grazie alle proprietà della trasformata di Fourier e della convoluzione abbiamo che, 
il segnale ricostruito ideale $x_r \left(t\right)$ è il seguente: 

{
    \Large 
    \begin{equation}
        x_r \left(t\right) 
        = 
        \sum_{n = -\infty}^{+ \infty} 
        x \left[  n \right]
        \dfrac{\sin \left( \dfrac{\pi \left(t - nT\right)}{T}\right)}{\dfrac{\pi \left(t - nT\right)}{T}}
    \end{equation}
}

$x_r \left(t\right)$ avrà questo andamento: 

\begin{figure}[h]
    \centering
    \includegraphics[scale = 0.4]{Segnale ricostruito dal filtro interpolatore.png}
\end{figure}  

dove, in figura: 

\begin{itemize}
    \item in blu è segnato l'anamento del segnale di $x_r \left(t\right)$
    \item in rosso è segnato l'andamento dei sinc del filtro interpolatore centrati in multipli di T (nella formula di $x_r \left(t\right)$ si noti il nT)
\end{itemize}

Dopo tutte queste considerazioni, possiamo dire che, se si svolge un campionamento ideale, avremo che: 

{
    \Large 
    \begin{equation}
        x_r \left(t\right)
        = 
        x_c \left(t\right)
    \end{equation}
}

cioè il segnale ricostruito $x_r \left(t\right)$ è uguale al segnale campionato $x_c \left(t \right)$. \newline  

Ciò può avvenire se non si ha aliasing, cioè: 

\begin{itemize}
    \item se il campionamento e il filtro interpolatore passano per i punti originali 
    \item se non c'è ambiguità 
\end{itemize}

\newpage 

\section{Riepilogo: campionamento e ricostruzione ideali} 

\footnote{ 
    Slide del prof | 4 - Campionamento | pag 9 - 10\\
    Appunti| 2025-10-14 | pag 2
}

Di seguito le formule riassuntive per un campionamento ideale. \newline 

Se consideriamo un campionamento ideale: 

\begin{figure}[h]
    \centering
    \includegraphics[scale = 0.3]{Campionamento ideale schema a blocchi.png}
\end{figure}  

dove: 

\begin{itemize}
    \item $x_c \left(t\right)$ è il segnale analogico da campionare 
    \item $x \left[ n\right]$ è il segnale campionato
\end{itemize}

Sapendo che: 

{
    \Large 
    \begin{equation}
        x \left[n\right] = x_c \left(n T\right)
    \end{equation}
}

cioè il segnale campionato $x \left[n\right]$ è uguale al segnale analogico $x_c (t)$ 
se viene campionato ogni periodo T, cioè a n multipli di T. \newline 

Allora, 
se consideriamo $X \left(e^{\jmath \omega}\right)$ la trasformata di Fourier del segnale campionato $x \left[n\right]$ 
e e $X \left(\jmath \Omega\right)$ la trasformata di Fourier del segnale analogico da campionare $x_c (t)$, 
avremo che la loro relazione è la seguente: 

{
    \Large 
    \begin{equation}
        X \left(e^{\jmath \omega}\right) 
        = 
        \dfrac{1}{T} 
        \sum_{k = -\infty}^{+ \infty}
        X_c \left(\jmath \dfrac{\omega}{T} - \jmath \dfrac{2 \pi k}{T}\right)
    \end{equation}
}

Se non c'è aliasing, possiamo scrivere che: 

{
    \Large 
    \begin{equation}
        X_c \left(\jmath \Omega\right) = 0 
        \text{ per }
        \abs{\Omega} > \Omega_N 
        \text{ e }
        \Omega_N < \dfrac{\Omega_s}{2}
    \end{equation}
}

dove: 

\begin{itemize}
    \item $\Omega$ sono le pulsazioni nello spettro continuo 
    \item $\Omega_N$ è la pulsazione di Nyquist 
    \item $\Omega_s$ è la pulsazione di campionamento
\end{itemize}

Allora, possiamo scrivere che: 

{
    \Large 
    \begin{equation}
        \begin{split}
            x \left[n \right] &= x_c \left(n T\right)
            \\
            &\updownarrow
            \\
            X \left(e^{\jmath \omega}\right)
            &= 
            \dfrac{1}{T}
            X_c \left(\dfrac{\jmath \omega}{T}\right)
            \text{ per }
            \abs{\omega} < \pi, k = 0
        \end{split}
    \end{equation}
}

\begin{tcolorbox}
    Tencicamente nello spettro di $X \left(e^{\jmath \omega}\right)$ ci dovrebbero essere le replice di $X_c \left(\dfrac{\jmath \omega}{T}\right)$, 
    ma, dal punto di vista informativo, è inutile indicare le replice dello spettro: ecco perchè ci poniamo nell'intervallo $\abs{\omega} < \pi$
\end{tcolorbox}

\newpage 

Invece per quanto riguarda la ricostruzione del segnale originale: 

\begin{figure}[h]
    \centering
    \includegraphics[scale = 0.3]{Interpolatore ideale schema.png }
\end{figure}  

dove: 

\begin{itemize}
    \item $x \left[ n\right]$ è il segnale campionato 
    \item $x_r \left(t\right)$ è il segnato ricostruito 
\end{itemize} 

Se non c'è aliasing, avremo che: 

{
    \Large 
    \begin{equation}
        x_r \left(t\right)
        = 
        x_c \left(t\right)
    \end{equation}
}

in cui, dal punto di vista matematico, $x_r \left(t\right)$ ha questo andamento: 

{
    \Large 
    \begin{equation}
        x_r \left(t\right) 
        = 
        \sum_{n = -\infty}^{+ \infty} 
        x \left[  n \right]
        \dfrac{\sin \left( \dfrac{\pi \left(t - nT\right)}{T}\right)}{\dfrac{\pi \left(t - nT\right)}{T}}
    \end{equation}
} 

che in frequenza sarà: 

{
    \Large 
    \begin{equation}
        X_r \left(\jmath \Omega\right)
        = 
        \begin{cases}
            T \cdot X \left(e^{\jmath \Omega T}\right)
            \text{ per }
            \abs{\Omega} < \dfrac{\Omega_s}{2}
            \\
            0 \text{ altrove}
        \end{cases}
    \end{equation}
}

dove: 

\begin{itemize}
    \item T moltiplicato per $X \left(e^{\jmath \Omega T}\right)$ indica il guadagno (o la perdita) del filtro interpolatore 
    \item $\Omega_s$ è la pulsazione di cutoff del filtro 
\end{itemize}

\begin{tcolorbox}[colback=red!5!white,colframe=red!75!black] \label{eq: 6}
    La lettera T nella formula $T \cdot X \left(e^{\jmath \Omega T}\right)$ intendono due cose differenti?
\end{tcolorbox} 

\newpage 

\section{Elaborazione segnali analogici con sistemi discreti} 

\footnote{ 
    Slide del prof | 4 - Campionamento | pag 11 - 12\\
    Appunti| 2025-10-14 | pag 2
}

Se consideriamo tutti i "blocchi" necessari per elaborare segnali analogici con un sistema discreto, 
possiamo considerare i seguenti schemi a blocchi: 

\begin{figure}[h]
    \centering
    \includegraphics[scale = 0.3]{Sistema a blocchi dell'elaborazione di un segnale analogico con sistema discreto.png}
\end{figure}  

dove: 

\begin{itemize}
    \item $x_c \left(t\right)$ e $y_r \left(t\right)$ sono segnali analogici (cioè grandezze fisiche tangibili come ad esempio una pressione di un tubo d'acqua, una tensione, una corrente, ecc...) 
    \item $x \left[n \right]$, $h \left[n \right]$, $y \left[n \right]$ sono delle sequenze (cioè sequenze che sono comprensibili ad un micro-controllore e/o ad un sistema digitale e che si trovano nel mondo "astratto" del DSP) 
\end{itemize}

\begin{tcolorbox}
    L'enorme vantaggio che si ha con questo tipo di schema a blocchi è che l'elaborazione del segnali analogici viene fatto in digitale con un software. \newline 

    Questa è una classica pipeline in quasi tutti i sistemi digitali moderni. \newline 

    Ad esempio, la Apple piuttosto che cambiare i filtri all'interno delle AirPods pro, 
    esegue un aggiornamento del software che permette di cambiare l'elaborazione dei segnali fisici, 
    come la cancellazione del rumore. \newline 

    Lo stesso esempio, nel mondo analogico, sarebbe stato diverso: Apple, per migliorare il prodotto, dovrebbe manualmente e fisicamente prendere le Airpods Pro di tutti i clienti, cambiargli fisicamente il filtro della cancellazione del rumore e poi spedirlo ai clienti di nuovo. \newline 

    Ecco perchè è molto più vantaggioso eseguire una elaborazione dei segnali analogici con sistemi discreti. \newline 

    Oppure, nel caso moderno, anche nelle macchine ciò accade, con gli aggiornamenti chiamati "over-the-air", come viene spiegato in questo video: \newline 

    What is an automotive over-the-air update? by Autovista \\
    \url{https://youtu.be/18NCdr9_fuA?si=jBNFzJnGcIm1sZsy}
\end{tcolorbox}

Analizzando in frequenza questi blocchi, 
avremo che il campionatore (se non c'è aliasing) avrò il seguente andamento: 

{
    \Large 
    \begin{equation}
        X \left(e^{\jmath \omega}\right) = \dfrac{1}{T} X_c \left(\dfrac{\jmath \omega}{T}\right)
        \text{ per } \abs{\omega} < \pi
    \end{equation}

}

dove i valori della pulsazione $\omega$ e del periodo $T$ sono scelti opportunamente secondo il teorema del campionamento. \newline 

Il sistema LTI (Linear Time Invariant) ha il seguente andamento in frequenza: 

{
    \Large 
    \begin{equation}
        Y \left(e^{\jmath \omega}\right)
        =
        H \left(e^{\jmath \omega}\right) \cdot X \left(e^{\jmath \omega}\right)
        \text{ per } \abs{\omega} < \pi
    \end{equation}
}

Invece il ricostruttore (indicato in figura come Interpolatore T) ha il seguente andamento in frequenza: 

{
    \Large 
    \begin{equation}
        Y_r \left(\jmath \Omega\right) 
        = 
        \begin{cases}
            T \cdot Y \left(e^{\jmath \Omega T}\right)
            \text{ per }
            \abs{\Omega} < \dfrac{\Omega_s}{2} 
            \\
            0 \text{ altrove}
        \end{cases}
    \end{equation}
}

Se invece consideriamo tutti i tre blocchi come un unico blocco, 
possiamo scrivere che in frequenza l'uscita del blocco avrà il seguente andamento: 

{
    \Large 
    \begin{equation}
        Y_r \left(\jmath \Omega\right) 
        = 
        \begin{cases}
            H \left(e^{\jmath \Omega T}\right) \cdot X_c \left(\jmath \Omega\right)
            \text{ per }
            \abs{\Omega} < \dfrac{\Omega_s}{2} 
            \\
            0 \text{ altrove}
        \end{cases}
    \end{equation}
}

dove: 

\begin{itemize}
    \item $H \left(e^{\jmath \Omega T}\right)$ è la risposta impulsiva di tutti i tre blocchi messi insieme nel dominio dei discreti in frequenza 
    \item $X_c \left(\jmath \Omega\right)$ è la trasformata di Fourier del segnale analogico iniziale $x_c \left(t\right)$
\end{itemize}

Se si considera un unico blocco di elaborazione, si può accettare aliasing se elimato da $H \left(e^{\jmath \omega}\right)$. \newline 

Quindi, se consideriamo la situazione ideale, 
fare i calcoli in frequenza è uguale che farli nel tempo e viceversa perchè è garantito dal teorema di Fourier e dal teorema del campionamento ideale, grazie al concetto di trasformata in cui il segnale è reversibile. \newline 

Invece, la pipeline reale sarà ben diversa. \newline 

\newpage 

\subsection{Conversione A/D reale} 

\footnote{ 
    Slide del prof | 4 - Campionamento | pag 13\\
    Appunti| 2025-10-14 | pag 2
}